% Part: applied-modal-logics
% Chapter: temporal-logic
% Section: introduction

\documentclass[../../../include/open-logic-section]{subfiles}

\begin{document}

\olfileid[pt]{aml}{tl}{int}

\olsection{Introdução}

As lógicas temporais tratam de afirmações sobre coisas que serão ou já
foram o caso. Arthur Prior é reconhecido como o criador da lógica
temporal, que chamou de \emph{lógica dos tempos verbais}. Nosso
tratamento da lógica temporal seguirá, em grande medida, a abordagem
modal original de Prior: introduzir operadores temporais no arcabouço
básico da lógica proposicional, que em geral trata as afirmações como
desprovidas de tempo verbal.

Por exemplo, na lógica proposicional, eu poderia falar de uma cadela,
Beezie, que às vezes fica sentada e às vezes não, como os cães costumam
fazer. Na lógica clássica, seria contraditório afirmar que Beezie está
sentada e também que Beezie não está sentada. É claro, porém, que ambas
as afirmações podem ser verdadeiras, apenas não ao mesmo tempo; o
acréscimo de operadores temporais à linguagem permite expressar isso
com relativa facilidade. Esses operadores também permitem explicar a
validade de inferências como a de ``Beezie receberá um petisco ou uma
bola" para ``Beezie receberá um petisco ou Beezie receberá uma bola." 

Entretanto, a lógica temporal suscita muitas questões filosóficas que
podem nos levar a preferir um arcabouço de lógica temporal a outro. Um
contingente futuro, por exemplo, é uma afirmação sobre o futuro que não
é necessária nem impossível. Se dissermos ``Richard irá ao mercado
amanhã", expressamos uma afirmação sobre algo que ainda não aconteceu e
cujo valor de verdade é discutível. Na realidade, é discutível até
mesmo se tal afirmação pode, em primeiro lugar, \emph{receber} um valor
de verdade. Se formos deterministas estritos, talvez aceitemos a ideia
de que essa sentença já é verdadeira ou falsa, mesmo antes do momento
em que o acontecimento em questão deveria ocorrer---pode ser apenas que
ainda não conheçamos seu valor de verdade. Em contraste, podemos crer
em um futuro genuinamente aberto, no qual os valores de verdade dos
contingentes futuros são indeterminados. 

Muitos desses compromissos acerca da estrutura e da natureza do tempo
acabam incorporados à escolha de modelos e arcabouços de lógica
temporal. Podemos perguntar, por exemplo, se devemos construir modelos
nos quais o tempo seja linear, ramificado ou até circular. Talvez seja
preciso decidir se nossos modelos temporais terão pontos inicial e
final e se o tempo será representado por instantes discretos ou como um
contínuo. 

\end{document}
